\noindent\underline{证明.}
设 $i:Y\hookrightarrow X$ 为包含映射. 有
\[
N_{Y/X}=H|_Y=i^*H.
\]
于是正合列
\[
0\longrightarrow T_Y
\longrightarrow i^*T_X
\longrightarrow i^*H
\longrightarrow0
\]
给出
\begin{equation}\tag{5.7.4}
C(Y)=i^*\!\left(\frac{C(X)}{1+H}\right).
\end{equation}
所以
\[
\chi(Y)=\int_Yi^*\!\left(\frac{C(X)}{1+H}\right),
\]
而投影公式给出
\begin{equation}\tag{5.7.5}
\chi(Y)
=\int_Xi_*i^*\!\left(\frac{C(X)}{1+H}\right)
=\int_X\frac{H\,C(X)}{1+H}.
\end{equation}
对 $\Delta$ 迭代此计算, 得到
\begin{equation}\tag{5.7.6}
\begin{aligned}
\chi(\Delta)
&=\int_Y\frac{H\,C(Y)}{1+H}\\
&=\int_Yi^*\!\left(
  \frac{H}{1+H}\cdot\frac{C(X)}{1+H}
  \right)
=\int_X\frac{H^2C(X)}{(1+H)^2}.
\end{aligned}
\end{equation}
于是
\[
\begin{aligned}
\chi(X)-2\chi(Y)+\chi(\Delta)
&=\int_XC(X)\left[
1-\frac{2H}{1+H}+\frac{H^2}{(1+H)^2}
\right]\\
&=\int_X\frac{C(X)}{(1+H)^2}.
\end{aligned}
\]
结合 (5.6.1), 即证得命题.

\section*{6. 一般底的情形}

\subsubsection*{6.1.}
设 $S$ 为概形, $\bar s$ 为 $S$ 的几何点, 由某点 $s\in S$ 的剩余域 $k(s)$ 的一个代数闭包定义. 设 $\mathbb P^r_S$ 为 $S$ 上的 $r$ 维射影空间, $\check{\mathbb P}^{r}_S$ 为其对偶射影空间, 并设
\[
I\subset\mathbb P^r_S\times_S\check{\mathbb P}^{r}_S
\]
为关联簇. 设 $D\subset\check{\mathbb P}^{r}_S$ 为由相应 Grassmann 簇的一个 $S$-值点定义的直线, $\Delta$ 为相应超平面束的轴(参见 2.1).

设 $X$ 为纯维数 $n\geq1$ 的固有光滑 $S$-概形, 配有射影嵌入 $X\hookrightarrow\mathbb P^r_S$. 置
\[
\widetilde X=(X\times_SD)\cap I.
\]
故 $\widetilde X$ 为 $D$ 上的概形, 其纤维为由 $D$ 参数化的 $X$ 的超平面截面; 设 $u:\widetilde X\to D$ 为投影.

以下标 $\bar s$ 表示沿 $\bar s\to S$ 的基变换. 假设 $\Delta_{\bar s}$ 与 $X_{\bar s}$ 横截. 则在 $s$ 于 $S$ 中的一个邻域 $U$ 上, $\Delta$ 与 $X$ 横截. 以下将 $S$ 换成此邻域, 并写作 $U=S$. 此时概形 $\widetilde X$ 在 $S$ 上光滑, 且 $\Delta$ 的逆像在 $\widetilde X$ 中的一个邻域在 $D$ 上光滑.

再假设属于束 $D_{\bar s}$ 的 $X_{\bar s}$ 的某个超平面截面光滑. 再将 $S$ 换成适当的 $s$ 的邻域后, 由 Jacobian 理想 $J^{n-1}(\widetilde X/D)$ 定义的闭子概形 $F\subset\widetilde X$ 在 $S$ 上有限且平坦.

\subsubsection*{6.2.}
束 $D$ 的\underline{例外轨迹}为 $D/S$ 的相对 Cartier 除子 $E$, 由有限平坦 $\mathcal O_D$-模 $u_*\mathcal O_F$ 定义. 在 $D$ 上局部地, 若
\[
\mathcal O_D^m\xrightarrow{\ v\ }\mathcal O_D^m
\longrightarrow u_*\mathcal O_F\longrightarrow0
\]
为 $u_*\mathcal O_F$ 的一个表示, 则 $E$ 的方程为 $\det(v)=0$. $E$ 的形成与任意基变换交换.

$D_{\bar s}$ 上的除子 $E_{\bar s}$ 可写为
\[
E_{\bar s}=\sum_P n(P)[P],
\]
其中 $n(P)$ 为 Jacobian 理想 $J^{n-1}(\widetilde X/D)$ 在 $u$ 的映到 $P$ 的各临界点处余长度之和.

\subsubsection*{命题 6.3.}
\underline{假设 $D_{\bar s}$ 为 $X_{\bar s}$ 的超平面截面的 Lefschetz 线束, 且 $\operatorname{char}k(s)\ne2$ 或 $n$ 为偶数. 则 $s$ 在 $S$ 中有一个邻域 $U$, 使:}
\begin{enumerate}
\item[(a)] \underline{对 $U$ 的每个几何点 $\bar t$, $D_{\bar t}$ 为 $X_{\bar t}$ 的超平面截面的 Lefschetz 线束;}
\item[(b)] \underline{在 $U$ 上, 例外轨迹 $E$ 在 $S$ 上平展.}
\end{enumerate}
为证明 (a) 和 (b), 只需注意在这些假设下, $D_{\bar s}$ 满足条件 2.2(c), 当且仅当除子 $E_{\bar s}$ 约化. 引用 3.3, 则只需使用从左到右这一较容易的蕴含.

\subsubsection*{6.3 (例外情形).}
\footnote{印本在命题 6.3 后重复使用编号 6.3.}
此命题在刚才提到的例外情形中不成立. 不过, 若 $S$ 的特征为二且 $n$ 为奇数, 则存在 $s$ 的一个邻域 $U$, 使 (a) 在其上成立, 并可定义一个在 $U$ 上平展的除子 $E'$, 满足 $E=2E'$. 我们不会使用此结果.

\subsubsection*{参考文献}

\noindent[1] A. Grothendieck, \emph{La théorie des classes de Chern},
Bull. Soc. Math. France \textbf{86} (1958), 137--154.

